% RESUMO - PT
\begin{resumo}
  \vspace{-15pt}
  
No âmbito da automação industrial, a integração entre sistemas de controle constitui uma necessidade fundamental para o monitoramento e controle eficiente de processos produtivos. Este trabalho apresenta o desenvolvimento de um sistema de supervisão e controle de nível baseado em lógica Fuzzy, integrado através do protocolo OPC (Open Platform Communications).

A lógica Fuzzy, técnica de inteligência artificial que possibilita o tratamento de incertezas e imprecisões inerentes aos processos industriais, demonstra particular adequação para o controle de plantas de nível. Tais sistemas, amplamente utilizados nas indústrias petroquímica, farmacêutica e de alimentos e bebidas, caracterizam-se por sua natureza não linear e frequentemente multivariável, apresentando desafios significativos para técnicas de controle convencionais.

O objetivo principal deste trabalho consiste na implementação de uma arquitetura de integração entre uma API responsável pelo controle Fuzzy e uma planta de nível simulada, utilizando o protocolo OPC como meio de comunicação. Esta abordagem visa demonstrar a viabilidade da aplicação de técnicas de inteligência artificial em ambientes industriais reais, aproveitando-se da padronização e ampla adoção do protocolo OPC na indústria moderna.


  \textbf{Palavras-chave:} Controle de Nível, Controlador Fuzzy, Lógica Fuzzy, C-Sharp, C\#, Automação Industrial, Python, Blazor, WPF, OPC, OPC UA, IHM, Integração de Sistemas.
\end{resumo}
